\section{Discussion}
\label{sec:discussion}

The preceding sections describe the architecture and tool ecosystem of \name in detail. Here we step back from per-component descriptions and examine five cross-cutting design tensions that shaped the system. Each subsection synthesizes insights that span multiple components and extracts transferable lessons for builders of similar agentic systems.

\subsection{Context Pressure as the Central Design Constraint}

Unlike CPU or memory, context is consumed by both the system (prompts, tool schemas, safety preambles) and the agent's own actions (tool outputs, conversation history). Every capability added to the system prompt and every tool result returned to the agent competes for the same finite budget. In our experience, tool outputs (file contents, command results, search hits) consume 70--80\% of the context in a typical session, dwarfing the system prompt and the agent's own reasoning. This makes context utilization the single most important metric for agent longevity, and it imposes a pervasive tension: richer tool outputs improve per-turn accuracy but shorten the session's useful life.

\begin{lessonbox}{Treat context as a budget, not a buffer}
Context reduction is not a binary operation. A graduated approach (monitoring utilization continuously, pruning stale tool outputs before they become irrelevant, masking old observations, and reserving LLM-based summarization for genuine overflow) dramatically outperforms the naive strategy of compacting everything once a hard limit is reached. The fast pruning pass described in \Cref{sec:adaptive_compaction} exemplifies this: by walking backwards through tool results and replacing outputs beyond the agent's working horizon with a \texttt{[pruned]} marker, it often reclaims enough space to avoid expensive LLM compaction entirely. Design graduated reduction stages rather than a single emergency compaction.
\end{lessonbox}

\begin{lessonbox}{Offload large outputs to the filesystem}
When a tool produces output exceeding a size threshold, writing the full content to a scratch file and returning only a short preview plus a file reference keeps the context focused on what the agent is actively using (\Cref{sec:tool_results}). The agent sees enough to decide whether the full output matters, and can read it on demand. This transforms a context-consumption problem into a retrieval problem, which is strictly cheaper: retrieval costs one tool call, while context consumption is paid on every subsequent LLM invocation for the rest of the session.
\end{lessonbox}

\paragraph{Prompt structure for caching.} For providers that support prompt caching, splitting the system prompt into a stable prefix and a dynamic suffix, marking the prefix with a cache control header, yields significant input cost savings over a multi-turn session (\Cref{sec:prompt_composition}). Since the system prompt is re-sent on every LLM call, caching the stable portion amortizes the cost of the system's richest instructions.

\paragraph{Dual-memory for thinking contexts.} When providing conversation context to a thinking model, separating compressed long-range context (an LLM summary of the full history) from detailed short-range context (the last several exchanges verbatim) keeps the thinking budget bounded regardless of conversation length while preserving both strategic goals and operational details (\Cref{sec:dual_memory}). A subtlety: iteratively summarizing a summary accumulates distortion over multiple rounds. Periodically regenerating the summary from the full conversation history, rather than compressing the previous summary, corrects this drift.

\begin{lessonbox}{Calibrate from API-reported token counts, not local estimates}
Using the API's reported \texttt{prompt\_tokens} from the previous call as the calibration anchor is essential for accurate context management. Providers inject invisible content (safety preambles, tool schema serializations, internal formatting) that makes local token counting systematically underestimate actual usage. In our system, the discrepancy was large enough to cause compaction to trigger too late, leading to context overflow errors (\Cref{sec:context_retrieval}). Always treat the provider's reported token count as ground truth; local estimates are useful only as a fallback when no prior API response is available.
\end{lessonbox}

\subsection{Steering Behavior Over Long Horizons}

System prompt influence decays as conversations grow. Instructions that reliably govern the agent's first few turns are routinely violated after 30 or more tool calls, when the instruction is far from the model's attention window and buried under dozens of tool results. Behavioral guidance is therefore a signal-to-noise engineering problem: how to maintain compliance without drowning the context in repeated directives.

\begin{lessonbox}{Inject reminders at the point of decision, not upfront}
Short, targeted reminders at maximum recency are more effective than long system prompt sections that the model has partially forgotten after 20 turns. Crucially, these reminders should use \texttt{role:~user} rather than \texttt{role:~system}: in our experiments, user-role reminders consistently produced stronger compliance, likely because the model assigns higher salience to recent user messages than to system context that has been pushed down by intervening turns (\Cref{sec:system_reminders}). However, reminder frequency must be capped per type; a reminder injected on every turn becomes background noise that the model learns to ignore, paradoxically reducing compliance.
\end{lessonbox}

\begin{lessonbox}{Separate thinking from action}
When tools are available, LLMs tend to act quickly rather than think deeply. Providing a separate thinking phase with no tool access produces substantially better reasoning traces, because the model is not pressured by the availability of action (\Cref{sec:react_executor}). This separation is more effective than asking the model to ``think carefully'' in the same call where tools are available. The mechanism matters: it is the absence of tool schemas from the API call, not an instruction to refrain from using them, that changes the model's behavior.
\end{lessonbox}

\paragraph{Explicit decision trees for tool selection.} Agents default to text search (grep) for nearly every lookup, even when a more precise tool exists. This wastes iterations and floods the context with false positives. Encoding a retrieval-tool decision tree directly into the agent's prompt (routing symbol names to semantic search, string patterns to text search, structural patterns to AST search, and file-name conventions to glob) reduced unnecessary grep calls and improved first-attempt retrieval accuracy. The key is that the decision criteria must be concrete and anchored to observable features of the query (``if the target is a function or class name, use \texttt{find\_symbol}''), not abstract (``use the most appropriate tool'').

\paragraph{Provider-conditional prompt sections.} Different LLM providers have meaningfully different capabilities: extended thinking, function calling conventions, context limits. Rather than cluttering the system prompt with provider-agnostic instructions, registering provider-specific sections that load only when the corresponding provider is active keeps the prompt budget focused (\Cref{sec:prompt_composition}). Unknown providers receive no section; graceful degradation is better than incorrect guidance.

\begin{lessonbox}{Separate agent construction from agent execution}
Long-running agents benefit from a clean separation between scaffolding (pre-runtime assembly) and the harness (runtime orchestration). Scaffolding (constructing the system prompt, building tool schemas, registering subagents) runs once before the first prompt and produces a fully initialized agent. The harness then wraps the reasoning loop with tool dispatch, context management, safety enforcement, and session persistence (\Cref{sec:agent_core}). This separation prevents construction-time concerns from tangling with runtime concerns: the harness never checks whether the agent is fully initialized, because eager construction guarantees it always is (\Cref{sec:agent_scaffolding}). The practical benefit is that each concern can evolve independently: adding a new tool requires only a registry change at construction time, while changing the compaction strategy requires only a harness change at runtime.
\end{lessonbox}

\subsection{Safety Through Architectural Constraints}

Runtime permission checks are the wrong primary abstraction for agent safety. A model that sees a dangerous tool in its schema can reason about invoking it, argue for why it should be allowed, and probe for edge cases in the permission logic. The more robust approach is to make violations structurally impossible: if write tools are absent from the agent's schema, the agent cannot attempt writes because it never sees a way to invoke them (\Cref{sec:agent_core}). This is the difference between a guard rail and a missing road: the model cannot reason about capabilities it does not know exist.

\begin{lessonbox}{Make unsafe tools invisible, not blocked}
Schema gating (removing tools from the agent's available set rather than checking permissions at call time) is fundamentally more robust than runtime permission checks. When a tool is absent from the schema, the agent cannot attempt to invoke it, argue for an exception, or probe for bypass conditions. Defense-in-depth with independently designed layers ensures that no single bypass compromises safety: schema gating prevents unauthorized attempts, the approval system intercepts operations requiring human review, file freshness validation prevents overwriting concurrent edits, and shadow git snapshots (\Cref{sec:undo}) make all filesystem changes reversible. These layers are deliberately independent; a bug in one does not weaken another.
\end{lessonbox}

\paragraph{Approval persistence prevents fatigue.} When a user marks an approval rule as ``always allow,'' persisting it to disk so it survives session restarts is critical. Without persistence, users must re-approve the same operations every session, causing approval fatigue that leads to blanket auto-approval, defeating the safety system entirely.

\paragraph{Lifecycle hooks as extensibility.} External scripts that observe or intercept agent lifecycle events enable custom policies, logging, CI integration, and security enforcement without modifying agent code. The hook interface must be designed early: blocking vs.\ non-blocking semantics, input mutation support, and merged global/project-level configurations are all necessary for hooks to be useful rather than decorative.

\paragraph{Modal priority during interrupt.} When the user presses an interrupt key while a modal dialog (approval prompt, user question) is active, the system should cancel the dialog rather than interrupting the agent. Interrupting the agent while a dialog is pending creates orphaned UI state (dangling spinners, stale futures, unresolved promises) that requires manual cleanup.

\subsection{Designing for Approximate Outputs}

LLMs reliably produce approximately-correct outputs. Edit targets drift from actual file content: trailing whitespace, indentation differences, escape sequence variations. Recovery strategies reference tools the agent does not have. Search queries route to suboptimal tools. A system that demands exact correctness from the model will spend most of its time in error-recovery loops. The alternative is to design tools and interfaces that absorb LLM imprecision as a first-class property.

\begin{lessonbox}{Design tools to absorb LLM imprecision}
The edit operation is the clearest example. A strict exact-match edit tool causes the majority of agent errors in practice, not because the agent's intent is wrong, but because its reproduction of the target text is slightly off. Building a chain of progressively relaxed matchers, where each returns the \emph{actual substring found in the file} to preserve original formatting, converts these near-misses into successful edits (\Cref{sec:file_ops}). The chain short-circuits on first match, so exact matches incur zero overhead. The general principle: when the agent's intent is unambiguous but its literal output is imprecise, the tool should bridge the gap rather than reject the attempt.
\end{lessonbox}

\begin{lessonbox}{Adapt recovery hints to the agent's available tool set}
When truncating large output and suggesting recovery strategies, the system must check which tools the agent actually has (\Cref{sec:tool_results}). If the agent can delegate to a subagent, suggest that; if it cannot (because it is itself a subagent), suggest incremental processing with offset/limit parameters. Generic hints that recommend unavailable tools cause the agent to attempt impossible actions and enter error loops. The same principle applies to error messages: ``re-read the file and retry your edit with current content'' is vastly more actionable than ``try again,'' so classifying errors into specific categories and retrieving targeted recovery templates for each improves recovery rates (\Cref{sec:error_recovery}).
\end{lessonbox}

\paragraph{Auto-promoting server-like commands.} Development servers, build watchers, and test suites that run indefinitely will hit any foreground timeout. Detecting server-like commands via regex patterns and automatically promoting them to background execution with output capture prevents the agent from blocking on a process that was never meant to terminate (\Cref{sec:bash_tool}).

\paragraph{Auto-installing missing dependencies.} Tools that depend on external runtimes should auto-install the dependency on first use rather than failing with an opaque error. Check for the dependency, install if missing, and retry. This eliminates a common class of setup-related failures that confuse both the agent and the user, and is another instance of designing tools to absorb environmental imprecision.

\subsection{Lazy Loading and Bounded Growth}

Eager loading fails at scale. Loading all MCP tool schemas at startup consumed 40\% of the context budget before the agent processed its first user message. Loading all skill definitions populated the prompt with content the agent would never use in most sessions. The solution in both cases is lazy discovery: load metadata indexes at startup, defer full content to the point of use (\Cref{sec:mcp}).

For MCP tools, lazy discovery reduces the startup context cost from 40\% to under 5\%. The agent receives a compact summary of available servers and their capabilities; full tool schemas are loaded only when the agent selects a server for a specific task. For skills, a two-phase approach serves the same purpose: a metadata index (name, description, trigger conditions) is loaded at startup, while the full skill content (which may include multi-page prompt templates) is loaded only when the agent decides to invoke the skill.

External metadata (model capabilities, pricing, context limits) benefits from a stale-while-revalidate caching strategy (\Cref{sec:provider_cache}): on startup, serve from cache if fresh; if stale, serve the stale data and refresh in the background; if the refresh fails, continue with stale data. This guarantees offline startup and eliminates ``cannot connect'' failures that block the entire system.

\begin{lessonbox}{Bound every resource that grows with session length}
Unbounded resources fail in long sessions. Without explicit caps, undo history accumulates snapshots indefinitely, concurrent tool calls can overwhelm the system, and behavioral nudges can flood the context. The design rule is simple: every resource that grows with session length must have a cap. Iteration limits prevent runaway agent loops. Undo history is bounded to a fixed number of snapshots. Concurrent tool calls are capped to prevent resource exhaustion. Nudge budgets limit how many times each reminder type can fire. Read-only tool calls run in parallel while write calls serialize (\Cref{sec:conversation_lifecycle}).
\end{lessonbox}

\begin{lessonbox}{Prefer empirical threshold tuning over first-principles calculation}
The specific values for caps (context pressure thresholds, pruning protection windows, reminder frequencies, iteration limits) resist first-principles calculation. They depend on the interaction between model behavior, typical user workflows, and system overhead in ways that are difficult to predict. Starting with conservative values and tuning based on observed failure modes (sessions that compact too early, reminders that fire too often, loops that run too long) is more effective than attempting to derive optimal values analytically. In our system, the 70\% compaction trigger, 3 nudge attempts, and 6 thinking depth levels all emerged from iterative failure analysis.
\end{lessonbox}

\paragraph{Self-healing indexes.} Caching frequently-accessed metadata in a lightweight index file enables fast listing without scanning underlying data files (\Cref{sec:session_management}). If the index is missing or corrupted, automatically rebuilding it from the underlying data makes the index a performance optimization rather than a single point of failure. The same principle applies to any derived data structure: design it so that loss triggers regeneration, not failure.

\paragraph{Deterministic operations outside the agent loop.} Not everything needs to go through the agent. Session management, mode switching, model selection, and server configuration are deterministic operations that should be handled directly at the input boundary (via slash-prefix dispatch or equivalent) with no approval gates, no undo tracking, and no token cost (\Cref{sec:repl_commands}). Routing these through the LLM wastes context and introduces unnecessary nondeterminism.

\medskip
\noindent These five design tensions (context as a budget, behavioral steering over long horizons, safety through architectural constraints, designing for approximate outputs, and bounding unbounded growth) are not unique to \name. They reflect fundamental challenges in building any long-running agentic system. The following section surveys how the broader research community has approached these same challenges, positioning \name's design decisions within the evolving landscape of code intelligence, autonomous software engineering, and context engineering.
